\documentclass[11pt,a4paper]{article}
\input{../../sablony/pripravy/priprava.sty}
\usepackage{tikz,pgfplots}
\pgfplotsset{compat=1.18}
\nastavPripravu{1. ročník}{Kvadratické rovnice a nerovnice}{104}{Neúplná kvadratická rovnice II}
\begin{document}
\blok{Neúplná kvadratická rovnice II}
Vytýkání neznámé a pravidlo nulového součinu; porovnání s odmocněním.
\begin{definiceBox}[title={Věta - Pravidlo nulového součinu}]
\textbf{Pravidlo nulového součinu:} $AB=0\iff A=0$ nebo $B=0$.
\end{definiceBox}
\blok{Metoda a návaznost na graf}
\begin{itemize}
\item Při $ax^2+bx=0$ vytkni $x$, nikdy neděl rovnou $x$: ztratil bys kořen $x=0$.
\item V každém příkladu nejprve rozhodni, zda je jednodušší odmocnění či vytýkání.
\end{itemize}
\blok{Společně řešené příklady}
\begin{prikladReseniBox}[title={Příklad 1}]
$x^{2} - 7x=0$
\tcblower
\[
\begin{aligned}
x^{2} - 7x&=x(x - 7)=0\\[2pt]
K&=\{0;7\}
\end{aligned}
\]
\end{prikladReseniBox}
\begin{prikladReseniBox}[title={Příklad 2}]
$3x^{2} + 6x=0$
\tcblower
\[
\begin{aligned}
3x^{2} + 6x&=x(3x + 6)=0\\[2pt]
K&=\{-2;0\}
\end{aligned}
\]
\end{prikladReseniBox}
\begin{prikladReseniBox}[title={Příklad 3}]
$18 - 2x^{2}=0$
\tcblower
\[
\begin{aligned}
x^2&=9\\|x|&=3\\[2pt]
K&=\{-3;3\}
\end{aligned}
\]
\end{prikladReseniBox}
\blok{Samostatné procvičování}
\begin{zadaniReseniBox}[title={Samostatná práce 1}]
$2x^{2} - 5x=0$
\tcblower
\[
\begin{aligned}
2x^{2} - 5x&=x(2x - 5)=0\\[2pt]
K&=\{0;\frac{5}{2}\}
\end{aligned}
\]
\end{zadaniReseniBox}
\begin{zadaniReseniBox}[title={Samostatná práce 2}]
$- 4x^{2} + 8x=0$
\tcblower
\[
\begin{aligned}
- 4x^{2} + 8x&=x(8 - 4x)=0\\[2pt]
K&=\{0;2\}
\end{aligned}
\]
\end{zadaniReseniBox}
\blok{Volitelná rezerva}
\poznamka{Studenti si vymyslí obdobný příklad se změnou jednoho koeficientu a vysvětlí, jak se změní počet kořenů či řešení.}
\blok{Rozšiřující příklady ze sbírky}
\poznamka{Nejdříve vytkni x, neztrať nulový kořen.}

\begin{zadaniReseniBox}[title={Rozšíření 1}]
Řešte v $\mathbb R$: $x \cdot \left(4x + 1\right)=0$.
\tcblower
\[
\begin{aligned}4x^2+x&=0\\x(4x+1)&=0\\K&=\{-\frac14;0\}\end{aligned}
\]
\end{zadaniReseniBox}

\blok{Procvičení před přechodem na úplné rovnice}
\begin{zadaniReseniBox}[title={Rozšíření 2}]
$5x^2-15x=0$.
\tcblower
\[\begin{aligned}
5x(x-3)&=0\\K&=\{0;3\}
\end{aligned}\]
\end{zadaniReseniBox}
\blok{Domácí úkol}
\begin{zadaniReseniBox}[title={Domácí úkol}]
\begin{enumerate}[label=\alph*),itemsep=2pt]
\item $x^{2} + 5x=0$
\item $3x^{2} - 12x=0$
\item $2x^{2} - 50=0$
\item $- x^{2} - 6x=0$
\item $5x^{2} + 45=0$
\item $4x^{2} + 7x=0$
\item Řešte v $\mathbb R$: $x \cdot \left(4x + 1\right)=0$.
\end{enumerate}
\tcblower
\begin{enumerate}[label=\alph*),itemsep=2pt]
\item
\[\begin{aligned}
x^{2} + 5x&=x(x + 5)=0\\[2pt]
K&=\{-5;0\}
\end{aligned}\]
\item
\[\begin{aligned}
3x^{2} - 12x&=x(3x - 12)=0\\[2pt]
K&=\{0;4\}
\end{aligned}\]
\item
\[\begin{aligned}
x^2&=25\\|x|&=5\\[2pt]
K&=\{-5;5\}
\end{aligned}\]
\item
\[\begin{aligned}
- x^{2} - 6x&=x(- x - 6)=0\\[2pt]
K&=\{-6;0\}
\end{aligned}\]
\item
\[\begin{aligned}
x^2&=-9\\[2pt]
K&=\varnothing
\end{aligned}\]
\item
\[\begin{aligned}
4x^{2} + 7x&=x(4x + 7)=0\\[2pt]
K&=\{- \frac{7}{4};0\}
\end{aligned}\]
\item \[ \begin{aligned}4x^2+x&=0\\x(4x+1)&=0\\K&=\{-\frac14;0\}\end{aligned}\]
\end{enumerate}
\end{zadaniReseniBox}
\poznamka{Orientační čas: úvod 5 min, společné příklady 18 min, samostatně 15 min, kontrola a shrnutí 7 min. Rezerva je dobrovolná. Číselné parametry modelových aplikací nejsou skutečně naměřená data.}
\end{document}
